% Baccalauréat malgache, série C, 2010 : énoncé.
\TCDBacAttribution{C}{2010}
\providecommand{\dx}{\mathrm{d}x}
\providecommand{\Re}[1]{\mathrm{Re}\left( #1 \right)}
\providecommand{\Im}[1]{\mathrm{Im}\left( #1 \right)}
\providecommand{\dt}{\mathrm{d}t}
\providecommand{\Ouv}{(O;\vec{u},\vec{v})}
\providecommand{\mes}[2]{\mathrm{mes}\left( \widehat{ \overrightarrow{#1} ; \overrightarrow{#2} } \right)}
\providecommand{\D}{(\mathscr{D})}
\providecommand{\Courbe}{(\mathscr{C})}
\providecommand{\congru}[3]{#1 \equiv #2 \quad \left[ \mathrm{mod}~#3\right]}
%
%%%%%%%%%%%%%%%%%%%%%%%%%%%%%%%%%%%%%%%%%%%%%%%%%%%%%%%%%%%%%%%%%%%%%%%%%%%%%%%%%%%%%%%%%%%%%%%%%%%%%%%%%%%%%%%%%%%%%%%%%%%%%%%%
\noindent \underline{\textsc{\textbf{Exercice}}} (4 points) \\%[-1\baselineskip]
      \underline{A- Arithm\'etique.} 
            \begin{enumerate}
                  \item Effectuer la division euclidienne de $-283$ par $19$.
                  \item \begin{enumerate}
                              \item Traduire en terme de congruence la division euclidienne de $11$ par $7$.
                              \item Prouver que, pour tout entier naturel $n$, le nombre 
                                    $(2^{3n+1} + 3^{6n} + 11^{3n+1})$ est divisible par $7$.
                        \end{enumerate}
                  \item Dans le syst\`eme d\'ecimal, un entier naturel non nul $X$ s'\'ecrit : 
                        $X = \overline{a_p a_{p-1} \cdots a_2 a_1 a_0}$ ; avec $a_p \neq 0$.
                        Justifier qu'on a : $\congru{X}{0}{11}$ si 
                        $
                              \displaystyle{
                                    \congru{\sum_{i=0}^{p} (-1)^i a_i }{0}{11}
                              }
                        $.
            \end{enumerate}
      \underline{B- Probabilit\'e.} \newline
            Dans un restaurant, un menu est compos\'e d'une entr\'ee, d'un plat et d'un dessert.
            On a le choix entre 3 entr\'ees, 5 plats et 4 desserts.
            Une jeune fille va dans ce restaurant et elle choisit un menu au hasard.
            \begin{enumerate}
                  \item Parmi les plats se trouve du poisson frit.
                        Quelle est la probabilit\'e pour qu'elle mange du poisson frit\,?
                  \item Dans les entr\'ees, il existe deux crudit\'es diff\'erentes et 
                        la glace aux fraises fait partie du dessert.
                        Quelle est la probabilit\'e pour qu'elle mange de la crudit\'e et de la glace aux fraises\,?
                  \item Quelle est la probabilit\'e pour qu'elle mange de la crudit\'e, 
                        du poisson frit et de la glace aux fraises\,?
            \end{enumerate}
%%%%%%%%%%%%%%%%%%%%%%%%%%%%%%%%%%%%%%%%%%%%%%%%%%%%%%%%%%%%%%%%%%%%%%%%%%%%%%%%%%%%%%%%%%%%%%%%%%%%%%%%%%%%%%%%%%%%%%%%%%%%%%%%
\noindent \underline{\textsc{\textbf{PROBLEME I}}} (7 points) 

      Le plan est rapport\'e \`a un rep\`ere orthonorm\'e direct. 
      
      \underline{\textbf{Partie A}} \\
      On consid\`ere un triangle \'equilat\'eral direct $TED$.
      $r_1$ est la rotation de centre $E$ et d'angle $\frac{2\pi}{3}$,
      $r_2$ la rotation de centre $D$ et d'angle $\frac{2\pi}{3}$.
      Utiliser le th\'eor\`eme de la d\'ecomposition pour d\'eterminer le centre de la rotation $r_1 \circ r_2$.
      
      \underline{\textbf{Partie B}} \\
      $OIJ$ est un triangle \'equilat\'eral direct.
      \begin{enumerate}
            \item Construire, \`a l'ext\'erieur de $OIJ$, le parall\'elogramme direct $OJKL$ et 
                  le triangle \'equilat\'eral direct $OLP$.
            \item On appelle $r$ la rotation de centre $O$ et d'angle $\frac{\pi}{3}$, 
                  et $t$ la translation de vecteur $\overrightarrow{JO}$.
                  \begin{enumerate}
                        \item Pr\'eciser la nature de la transformation $g$ d\'efinie par $g = r \circ t$.
                        \item D\'eterminer les images par $g$ des points $I$ et $K$.
                              En d\'eduire les \'el\'ements g\'eom\'etriques de $g$, ainsi que la nature du triangle $IKP$.
                  \end{enumerate}
      \end{enumerate}
      
      \underline{\textbf{Partie C}} \\
      \begin{description}
            \item[I.] On consid\`ere $M$, $E$ et $N$ trois points non align\'es.
                      Soit $H$ le barycentre du syst\`eme $\{ (M,3);(E,1);(N,1) \}$,
                      $Q$ celui de $\{ (M,3);(N,1) \}$ et $R$ celui de $\{ (M,3);(E,1) \}$.
                      \begin{enumerate}
                            \item D\'emontrer que les droites $(EQ)$ et $(NR)$ passent par $H$.
                            \item Soit $P$ le milieu du segment $[EN]$.
                                  \begin{enumerate}
                                        \item Prouver que $M$, $P$ et $H$ sont align\'es.
                                        \item Exprimer $\overrightarrow{PH}$ en fonction de $\overrightarrow{PM}$.
                                  \end{enumerate}
                      \end{enumerate}
            \item[II.] On donne un triangle rectangle direct $ABC$ et isoc\`ele en $A$ tel que : $AB = AC = a$ avec $a>0$.
                  \begin{enumerate}
                        \item D\'eterminer le point $G$ barycentre du syst\`eme 
                              $\{ (A,4);(B,-1);(C,-1) \}$. Construire $G$.
                        \item D\'eterminer l'ensemble $(\mathscr{E})$ des points $M$ du plan tels que : 
                              $4 \overline{MA}^2 - \overline{MB}^2 - \overline{MC}^2 = 2 a^2$.
                              Construire $(\mathscr{E})$.
                  \end{enumerate}
      \end{description}
%%%%%%%%%%%%%%%%%%%%%%%%%%%%%%%%%%%%%%%%%%%%%%%%%%%%%%%%%%%%%%%%%%%%%%%%%%%%%%%%%%%%%%%%%%%%%%%%%%%%%%%%%%%%%%%%%%%%%%%%%%%%%%%%
\noindent \underline{\textsc{\textbf{PROBLEME II}}} (9 points) \newline
      \underline{\textbf{PARTIE I}}
      Soit $f$ la fonction num\'erique d\'efinie par : 
      $
      \begin{multiform}
            f(0) = 0 \\
            f(x) = \frac{x-1}{\ln (x)} \text{ si } 0 < x < 1 \\
            f(1) = 1
      \end{multiform}
      $.
      $\Courbe$ d\'esigne sa courbe repr\'esentative dans un rep\`ere orthonorm\'e $\mathscr{R} = (O;\ivec,\jvec)$
      d'unit\'e graphique $10\cm$.
      \begin{description}
            \item[A. ]
                  \begin{enumerate}
                        \item Donner l'ensemble de d\'efinition de $f$.
                        \item Prouver que $f$ est continue en $0$ et en $1$.
                        \item Justifier que $f$ est d\'erivable sur $]0;1[$.
                              D\'emontrer que $f'(x)$ a le m\^eme signe que $H(x)$ o\`u $H$ est la fonction d\'efinie sur 
                              $]0;1[$ par : $H(x) = \frac{1}{x} - 1 + \ln (x)$.
                        \item Etudier les variations et le signe de $H$.
                              V\'erifier que $f$ est strictement croissante.
                        \item Etudier la d\'erivabilit\'e de $f$ en $0$.
                  \end{enumerate}
            \item[B. ]
                  \begin{enumerate}
                        \item Prouver que pour tout r\'eel $a$ de $\left[ 0 ; \frac{1}{2} \right[$ on a : 
                              $
                              0 \leq \frac{1}{1-a} - (1+a) \leq 2a^2
                              $
                              et
                              $
                              0 \leq - \ln (1-a) - \left( a + \frac{a^2}{2} \right) \leq \frac{2a^3}{3} \;.
                              $
                        \item On consid\`ere la fonction num\'erique $\Psi$ d\'efinie sur $]0;1]$ par :
                              $\Psi (t) = \frac{1}{f(t)}$.
                              Prouver que, pour tout r\'eel $b$ de $\left[ -\frac{1}{2} ; 0 \right]$, on a :
                              $0 \leq \Psi (1+b) - \Psi (1) + \frac{b}{2} \leq \frac{2b^3}{3}$.
                              En d\'eduire que $\Psi$ est d\'erivable en $1$ et donner $\Psi' (1)$.
                        \item D\'eduire de ce qui pr\'ec\`ede que $f$ est d\'erivable en $1$ et que $f'(1) = \frac{1}{2}$.
                        \item Tracer $\Courbe$ et ses tangentes \`a l'origine des axes et au point d'abscisse $1$.
                  \end{enumerate}
      \end{description}
      \underline{\textbf{PARTIE II}}
            Soit $(a_n)_{n\in\mathbb{N}^*}$ la suite num\'erique r\'eelle d\'efinie par : 
            $a_n = \frac{n^n}{n!}$ o\`u $n! = n \cdot (n-1) \cdots 3 \cdot 2 \cdot 1$.
            \begin{enumerate}
                  \item Etudier la monotonie de $(a_n)$.
                  \item Etablir les in\'egalit\'es : $2 \leq \left( 1 +\frac{1}{n} \right)^n \leq \expo{}$, $n\in\mathbb{N}^*$.
                        (On rappelle que pour tout r\'eel positif $x$, on a 
                        $1+nx \leq (1+x)^n$ et $\ln (1+x) \leq x$.)
                  \item D\'emontrer que, pour tout entier naturel $n$ sup\'erieur ou \'egal \`a 6, on a : 
                        $2^n \leq a_n \leq \expo{n}$.
            \end{enumerate}
%%%%%%%%%%%%%%%%%%%%%%%%%%%%%%%%%%%%%%%%%%%%%%%%%%%%%%%%%%%%%%%%%%%%%%%%%%%%%%%%%%%%%%%%%%%%%%%%%%%%%%%%%%%%%%%%%%%%%%%%%%%%%%%%





































